\documentclass{resume} % Use the custom resume.cls style

\usepackage[left=0.5in, top=0.5in, right=0.5in, bottom=0.5in]{geometry}
\usepackage{amsmath}
\newcommand{\tab}[1]{\hspace{.2667\textwidth}\rlap{#1}} 
\newcommand{\itab}[1]{\hspace{0em}\rlap{#1}}
\name{Anirban Chakraborty, Ph.D.}
%\address{+1(940) 594-8431 \\ 2945 S, Oleander Street, Cornelius, OR 97113}
\address{\href{mailto:anirban1.chakraborty@intel.com}{anirban1.chakraborty@intel.com} \\ \href{https://www.linkedin.com/in/anirbanchakraborty91/}{linkedin.com/in/anirbanchakraborty91}}

\begin{document}

\begin{rSection}{SUMMARY}
    \begin{itemize}
        \itemsep -3pt {}
        \item $1.5$ years experience in working as Resource Controller IP Engineer at Intel Corporation for the next generation of Server, Client, Networking and Graphics Roadmaps.
        \item Experienced Electrical Engineer with strong understanding of RTL Design, VLSI Simulation and Error Correction Coding.
        \item Knowledge on Physical Design, Place \& Route, IR drop, Physical Verification (DRC/LVS).
        \item $5+$ years of experience in System Designing using CAD tools like Synopsys Design Compiler, Synopsys PrimeTime, Synopsys VCS, Cadence Virtuoso, Cadence Innovus and Xilinx Vivado.
        \item $2+$ years of experience in Electrical Engineering data analysis using Python and R.
        \item An excellent team player with an ability to quickly learn new techniques and work in a fast-paced environment. 
    \end{itemize}
\end{rSection}

\begin{rSection}{TECHNICAL SKILLS}
    \begin{itemize}
        \itemsep -3pt {} 
        \item \textbf{HDL Languages:} SystemVerilog, VHDL
        \item \textbf{Design Techniques:} Synthesis, Physical Design, Place \& Route, IR drop, Physical Verification (DRC/LVS).
        \item \textbf{Developer and CAD Tools:} Cadence JasperGold, Synopsys Verdi, Synopsys Design Compiler, Synopsys PrimeTime, Synopsys VCS, Ansys PowerArtist, Cadence Virtuoso, Cadence Innovus, Xilinx Vivado, LTspice, PSpice, Electric Layout tool, Intel Quartus, ModelSim, MATLAB
        \item \textbf{Programming:} Python, R, C
    \end{itemize}
\end{rSection}

\begin{rSection}{EXPERIENCE}

    \textbf{IP Logic Design Engineer} \hfill June 2022 - Present \\
    \href{https://www.intel.com/content/www/us/en/homepage.html}{Intel Corporation} \hfill \textit{Hillsboro, OR, U.S.A.} \\
    \textbf{Resource Controller Development}
        \begin{itemize}
            \itemsep -3pt {} 
            \item Implemented Registers for the Resource controller IP used in the next generation Intel SOCs.
            \item Implemented the logic of component activation and level determination in the Power Management logics used in the next generation Intel SOCs.
            \item Implemented Formal Verification Properties for the Power Management Logics used in Intel SOCs.
            \item Performed thorough analysis and implementation of the stimulus for accurate workload simulation of the next generation Intel SOCs.
            \item Implemented exhaustive clock gating throughout the Resource controller IP to obtain a Dynamic Clock Gating Efficiency (DCGE) of $95\%$.
            \item Implemented triggers for debug in the Resource controller IP used in the next generation Intel SOCs.
            \item Worked on group wide upgrade of Cheetah Tools like SpyGlass Lint, Siemens Calibre, etc.
        \end{itemize}
        
    \textbf{Graduate Research and Teaching} \hfill Aug 2016 - May 2022 \\
    \href{https://www.unt.edu}{University of North Texas} \hfill \textit{Denton, TX, U.S.A.} \\
    \textbf{RTL Design --} \textit{Technology used} - 45 nm and 28 nm CMOS Process in Synopsys
        \begin{itemize}
            \itemsep -3pt {} 
            \item 4+ years of experience in VHDL and SystemVerilog.
            \item Brainstormed and collaborated on Hardware Architecture Constraints.
            \item Automated RTL Design Verification using TCL and Testbench Scripts and performed Synthesis on Synopsys Design Compiler.
            \item Worked on FPGA and ASIC platform.
            \item Taught VHDL and mentored in Designing and Testing Digital circuits in ASICs and FPGAs
        \end{itemize}
    
    \textbf{VLSI Design --} \textit{Technology used} - TSMC 180 nm CMOS Process in Cadence Virtuoso
        \begin{itemize}
            \itemsep -3pt {} 
            \item Investigated and verified the functionality of CMOS circuits.
            \item Performed Design, Testing and Simulation of Digital Circuits.
            \item Performed DRC check and Layouts of basic Digital Circuits
        \end{itemize}
    
    \textbf{Error Correction Coding --} Developing efficient Codes for Error Correction
        \begin{itemize}
            \itemsep -3pt {} 
            \item Experience in optimization of Codes.
            \item Investigated existing ECC and optimized their power and area requirement to improve the performance.
            \item Developed new and improved codes for NAND Flash architectures.
        \end{itemize}
        
    \textbf{Data Analysis --} Crowdsourced E-Gaming Framework for Untangled and Untangled III
        \begin{itemize}
            \itemsep -3pt {} 
            \item Conducted crowdsourcing event for Untangled and Untangled III to collect the data for investigating the problem of Application Mapping on Hardware.
            \item Converted the Mapped DFG into RTL Design using SystemVerilog and synthesized in ASIC platform.
        \end{itemize}
     
    \textbf{Guest Faculty} \hfill July 2015 - Dec 2015\\
    \href{https://www.seacomengineering.org}{Seacom Engineering College} \hfill \textit{Howrah, West Bengal, India}
        \begin{itemize}
            \itemsep -3pt {} 
            \item Instructor for the Electrical Machines Course.
        \end{itemize}
\end{rSection}

\begin{rSection}{Education}
    \textbf{Ph.D., Electrical Engineering --} GPA -- $3.812/4.00$ \hfill Aug 2016 - May 2022 \\
    \href{https://www.unt.edu}{University of North Texas} \hfill \textit{Denton, TX, U.S.A.}
        \begin{itemize}
            \itemsep -3pt {} 
            \item \textbf{Dissertation --} Novel Algorithms and Hardware Architectures for Computational Subsystems Used in Cryptography and Error Correction Coding.
        \end{itemize}
    
    \textbf{M.E., Electrical Engineering} \hfill July 2013 - April 2015 \\
    \href{https://www.iiests.ac.in}{Indian Institute of Engineering Science and Technology, Shibpur --} GPA -- $4.00/4.00$ \hfill \textit{Howrah, West Bengal, India}
        \begin{itemize}
            \itemsep -3pt {} 
            \item \textbf{Thesis --} Performance Analysis of a Single Axis Controlled Repulsive type Magnetic Bearing.
        \end{itemize}
    
    \textbf{B.Tech., Electrical Engineering --} GPA -- $8.63/10.00$ \hfill Aug 2009 - June 2013 \\
    \href{https://tint.edu.in}{Techno India College of Technology} \hfill \textit{Kolkata, West Bengal, India}
\end{rSection}

\begin{rSection}{SCHOLARSHIPS, HONORS \& NOMINATIONS}
    \begin{itemize}
        \itemsep -3pt {} 
        \item \textbf{Tuition Benefit Program --} \href{https://tgs.unt.edu/toulouse-graduate-school}{Toulouse Graduate School, UNT} \hfill Aug 2016 - May 2022
        \item \textbf{IEEE HKN Lambda-Zeta Student Poster Competition --} Third Place \hfill April 2022
        \item \textbf{Scholars Day Poster Judge --} \href{https://honors.unt.edu}{Honors College, UNT} \hfill April 2022
        \item \textbf{Three times Outstanding TA Nominations --} \href{https://electrical.engineering.unt.edu}{Department of Electrical Engineering, UNT} \hfill Aug 2019 - May 2022
        \item \textbf{Post Graduate Scholarship (GATE/GPAT) --} \href{https://www.aicte-india.org/schemes/students-development-schemes/PG-Scholarship-Scheme}{AICTE} \hfill July 2013 - April 2015
    \end{itemize}
\end{rSection}

\begin{rSection}{PUBLICATIONS}
    \textbf{Peer-Reviewed Journal Articles}
        \begin{itemize}
            \itemsep -3pt {} 
            \item \textbf{A. Chakraborty}, Z. Simpson, R. D. Balavendran Joseph, and G. Mehta, ``A Comparative Study of Gameplay of Different Sets of Players in an Engineering Mapping Game," \textbf{International Journal of Computer and Information Technology}(2279-0764), vol. 10, no. 4. Asian Online Journals, Jul. 31, \textbf{2021}. \href{https://doi.org/10.24203/ijcit.v10i4.127}{doi: 10.24203/ijcit.v10i4.127}.
            \item \textbf{A. Chakraborty}, E. Ojo, B. Quonoey, and G. Mehta, ``Improving Learning Experience of People with Cognitive Disabilities Using Serious Games: A Review," European Scientific Journal ESJ, vol. 17, no. 35. \textbf{European Scientific Institute, ESI}, Oct. 31, \textbf{2021}. \href{https://doi.org/10.19044/esj.2021.v17n35p1}{doi: 10.19044/esj.2021.v17n35p1}.
            \item \textbf{A. Chakraborty}, Z. P. Simpson, J. Tunks, and G. Mehta, ``Solving Electrical Engineering Puzzles Using Spatial Reasoning," \textbf{Journal of Education and Culture Studies}, vol. 2, no. 3. Scholink Co, Ltd., p. 227, Aug. 29, \textbf{2018}. \href{https://doi.org/10.22158/jecs.v2n3p227}{doi: 10.22158/jecs.v2n3p227}.
        \end{itemize}

    \textbf{Conference Proceedings}
        \begin{itemize}
            \itemsep -3pt {} 
            \item \textbf{A. Chakraborty}, M. Varanasi, O. N. Garcia, and G. Mehta, ``Area and Power analysis of a Scalable Primitive Polynomial computation circuit over the field GF(2)," \textbf{2022 IEEE 15th Dallas Circuit And System Conference (DCAS)}. IEEE, Jun. 17, \textbf{2022}. \href{https://doi.org/10.1109/DCAS53974.2022.9845656}{doi: 10.1109/dcas53974.2022.9845656}.
            \item Z. Simpson, \textbf{A. Chakraborty}, M. Varanasi, G. Mehta, and O. N. Garcia, ``Power, Performance, and Area Analysis of Hardware Design Techniques for GF(2k) Greatest Common Divisor computation," \textbf{2022 IEEE 15th Dallas Circuit And System Conference (DCAS)}. IEEE, Jun. 17, \textbf{2022}. \href{https://doi.org/10.1109/DCAS53974.2022.9845594}{doi: 10.1109/dcas53974.2022.9845594}.
            \item \textbf{A. Chakraborty}, T. Santra, D. Roy, and S. Yamada, ``Modelling and stability analysis of a repulsive type magnetic bearing," \textbf{2015 International Conference on Energy, Power and Environment: Towards Sustainable Growth (ICEPE)}. IEEE, Jun. \textbf{2015}. \href{https://doi.org/10.1109/EPETSG.2015.7510149}{doi: 10.1109/epetsg.2015.7510149}.
        \end{itemize}

    \textbf{Book Chapters}
        \begin{itemize}
            \itemsep -3pt {} 
            \item \textbf{A. Chakraborty}, ``Evolution of the Internet of Things and Fundamentals of Human-Machine Interaction," \textbf{Human-Machine Interaction and IoT Applications for a Smarter World}. CRC Press, pp. 3–14, Jun. 13, \textbf{2022}. \href{https://www.taylorfrancis.com/chapters/edit/10.1201/9781003268796-2/evolution-internet-things-fundamentals-human-machine-interaction-anirban-chakraborty}{doi: 10.1201/9781003268796-2}
        \end{itemize}
\end{rSection}

\end{document}